
\documentclass{article}
\usepackage{colortbl}
\usepackage{makecell}
\usepackage{multirow}
\usepackage{supertabular}

\begin{document}

\newcounter{utterance}

\centering \large Interaction Transcript for game `ifeval', experiment `ifeval', episode 339 with llama{-}3.1{-}8b{-}instruct.
\vspace{24pt}

{ \footnotesize  \setcounter{utterance}{1}
\setlength{\tabcolsep}{0pt}
\begin{supertabular}{c@{$\;$}|p{.15\linewidth}@{}p{.15\linewidth}p{.15\linewidth}p{.15\linewidth}p{.15\linewidth}p{.15\linewidth}}
    \# & \multicolumn{2}{c}{Player} && \multicolumn{2}{c}{Game Master} \\
    \hline

    \theutterance \stepcounter{utterance}  
    & & & \multicolumn{4}{p{0.6\linewidth}}{
        \cellcolor[rgb]{0.9,0.9,0.9}{
            \makecell[{{p{\linewidth}}}]{
                \texttt{\tiny{[P1$\langle$GM]}}
                \texttt{What is inside Shinto shrines? Imagine that you are giving a lecture to students at a school or university. Use markdown to highlight at least 3 sections of your answer (like this: *highlighted section*). Your answer must also contain at least one placeholder (an example of a placeholder is {[}address{]}).} \\
            }
        }
    }
     \\ \\

    \theutterance \stepcounter{utterance}  
    & \multicolumn{4}{p{0.6\linewidth}}{
        \cellcolor[rgb]{0.561,0.737,0.561}{
            \makecell[{{p{\linewidth}}}]{
                \texttt{\tiny{[P1$\rangle$GM]}}
                \texttt{**Introduction to Shinto Shrines**} \\
\\ 
\texttt{Good morning/afternoon students, today we're going to explore the fascinating world of Shinto shrines in Japan. Shinto shrines, also known as "jinja" in Japanese, are sacred spaces that hold great cultural and spiritual significance. As we delve into the world of Shinto shrines, let's take a closer look at what you can expect to find inside.} \\
\\ 
\texttt{\#\#\# *Sacred Objects and Artifacts*} \\
\\ 
\texttt{When you enter a Shinto shrine, you'll notice a variety of sacred objects and artifacts that are believed to possess spiritual power. Some of the most common items you'll see include:} \\
\\ 
\texttt{* **Shinboku** (Sacred Trees): These are trees that have been designated as sacred and are often found within the shrine grounds. They're believed to possess spiritual energy and are often used in Shinto rituals.} \\
\texttt{* **Shinmei** (Sacred Mirror): This is a sacred mirror that's believed to reflect the soul of the deity. It's often displayed in a special shrine room and is used in rituals to communicate with the deity.} \\
\texttt{* **Sakaki** (Sacred Tree Branches): These are branches from the sakaki tree, which is a sacred tree in Shintoism. They're often used in rituals and are believed to possess spiritual power.} \\
\\ 
\texttt{\#\#\# *Shrine Buildings and Structures*} \\
\\ 
\texttt{Shinto shrines are often composed of several buildings and structures, each with its own unique purpose. Some of the most common structures you'll see include:} \\
\\ 
\texttt{* **Honden** (Main Shrine Building): This is the main building of the shrine and is where the deity is enshrined. It's often a large, ornate building with a tiled roof.} \\
\texttt{* **Haiden** (Oratory): This is a building where visitors can pray and make offerings to the deity. It's often a smaller building with a thatched roof.} \\
\texttt{* **Heiden** (Offering Hall): This is a building where visitors can make offerings to the deity, such as food, incense, and other items.} \\
\\ 
\texttt{\#\#\# *Rituals and Practices*} \\
\\ 
\texttt{Shinto shrines are not just places of worship, but also centers of community and culture. Visitors can participate in various rituals and practices, such as:} \\
\\ 
\texttt{* **Matsuri** (Festivals): These are festivals that take place throughout the year, often to celebrate the deity or to mark important events.} \\
\texttt{* **Shinji** (Rituals): These are rituals that are performed by the shrine priests to communicate with the deity and to maintain the balance of nature.} \\
\texttt{* **O{-}mi{-}mairi** (Visiting the Shrine): This is a ritual where visitors come to the shrine to pray, make offerings, and seek blessings.} \\
\\ 
\texttt{**Conclusion**} \\
\\ 
\texttt{Shinto shrines are fascinating places that offer a glimpse into Japan's rich cultural and spiritual heritage. By understanding the sacred objects, shrine buildings, and rituals that take place within these shrines, we can gain a deeper appreciation for the importance of Shintoism in Japanese culture. If you're interested in learning more, I recommend visiting the {[}address{]} of the Meiji Shrine in Tokyo, one of the most famous Shinto shrines in Japan.} \\
            }
        }
    }
    & & \\ \\

    \theutterance \stepcounter{utterance}  
    & & & \multicolumn{2}{p{0.3\linewidth}}{
        \cellcolor[rgb]{0.9,0.9,0.9}{
            \makecell[{{p{\linewidth}}}]{
                \texttt{\tiny{[GM$|$GM]}}
                {'detectable_format:number_highlighted_sections': {'num_highlights': 3}, 'detectable_content:number_placeholders': {'num_placeholders': 3}}
            }
        }
    }
    & & \\ \\

    \theutterance \stepcounter{utterance}  
    & & & \multicolumn{2}{p{0.3\linewidth}}{
        \cellcolor[rgb]{0.9,0.9,0.9}{
            \makecell[{{p{\linewidth}}}]{
                \texttt{\tiny{[GM$|$GM]}}
                \texttt{game\_result = LOSE} \\
            }
        }
    }
    & & \\ \\

\end{supertabular}
}

\end{document}
